\documentclass[8pt,a4paper,twocolumn]{article}

\usepackage[spanish]{babel}
\usepackage[utf8]{inputenc}
\usepackage[T1]{fontenc}
\usepackage[left=0.75cm,right=0.75cm,top=0.55cm,bottom=0.7cm,columnsep=0.45cm]{geometry}
\usepackage{lmodern}
\usepackage{booktabs}
\usepackage{array}
\usepackage{tabularx}
\usepackage{hyperref}
\usepackage{enumitem}
\usepackage{xcolor}
\usepackage{titlesec}
\usepackage[most]{tcolorbox}

\hypersetup{
    colorlinks=true,
    linkcolor=blue,
    urlcolor=blue
}

\setlength{\parindent}{0pt}
\setlength{\parskip}{0.12em}
\setlength{\columnseprule}{0pt}
\pagestyle{plain}
\setlist[itemize]{leftmargin=1.0cm,itemsep=0.1em,topsep=0.15em}
\setlist[enumerate]{leftmargin=1.0cm,itemsep=0.1em,topsep=0.15em}
\setlist[description]{leftmargin=0cm,style=nextline,itemsep=0.1em,topsep=0.15em}
\titlespacing*{\section}{0pt}{0.35em}{0.2em}
\titlespacing*{\subsection}{0pt}{0.3em}{0.15em}

\newcommand{\cajaitem}[2]{
\noindent
\fcolorbox{black}{gray!10}{
\parbox{\dimexpr\linewidth-2\fboxsep-2\fboxrule\relax}{
\textbf{#1}\\#2
}}
}

\begin{document}

\vspace*{-1.0em}
\begin{tcolorbox}[colback=black!12,colframe=black!35,boxrule=0.3pt,arc=1pt,left=4pt,right=4pt,top=3pt,bottom=3pt]
{\bfseries\normalsize Material de estudio semanal\par}
{\scriptsize CIB12 - Aprendizaje de M\'aquina y Miner\'ia de Datos\par}
{\scriptsize Semana 4: Desarrollo de Aplicaciones Big Data\par}
{\scriptsize Fuente trabajada: Huawei (2022), Cap\'itulo 1\par}
\end{tcolorbox}
\vspace{0.15em}

\section*{2. Resumen de la semana}

Esta semana se estudia la parte del capítulo 1 enfocada en el \textbf{desarrollo de aplicaciones Big Data}. El contenido resalta que una buena aplicación no depende solo de programar, sino de combinar \textbf{base técnica} y \textbf{comprensión del servicio o negocio}. 

El material presenta dos ejes principales:

\begin{itemize}[leftmargin=*]
    \item \textbf{Habilidades requeridas}: programación en Java y Scala, manejo de SQL y operaciones en Linux, junto con comprensión del proceso de desarrollo y del contexto del servicio.
    \item \textbf{Proceso general de desarrollo}: desde la formulación del requerimiento hasta la visualización final de resultados.
\end{itemize}

La idea más importante es que el desarrollo de una aplicación Big Data sigue una \textbf{secuencia ordenada de etapas}: se identifican necesidades, se analizan requerimientos, se selecciona la tecnología, se valida su viabilidad, se preparan y procesan los datos, se almacenan los resultados y luego se muestran para uso operativo.

\section*{3. Conceptos clave}

\begin{table*}[t]
\small
\centering
\begin{tabularx}{\textwidth}{>{\raggedright\arraybackslash}p{0.28\textwidth} >{\raggedright\arraybackslash}X}
\toprule
\textbf{Concepto} & \textbf{Síntesis de estudio} \\
\midrule
Desarrollo de aplicaciones Big Data & Proceso de construir soluciones sobre datos masivos, integrando habilidades técnicas y entendimiento del negocio. \\
Programación básica & Base técnica necesaria para desarrollar aplicaciones; incluye Java, Scala, SQL y manejo de Linux. \\
Experiencia en desarrollo de servicios & Conocimiento del proceso de desarrollo y del contexto del servicio o problema que se desea resolver. \\
Requirement submission & Etapa inicial donde se negocia entre áreas y se analiza la factibilidad del proyecto. \\
Requirement analysis & Recolección y estudio de requerimientos mediante levantamiento de necesidades y estudio de mercado. \\
Technology selection & Selección de solución técnica y elaboración del diagrama de flujo de datos. \\
Feasibility analysis & Construcción de la plataforma técnica y pruebas para validar que la solución puede ejecutarse. \\
Indicator analysis & Conversión de necesidades del negocio en indicadores medibles. \\
Data interaction & Recolección, limpieza, almacenamiento de datos fuente y uso de middleware. \\
Data analysis & Implementación del código para procesar o analizar datos. \\
Result storage & Almacenamiento del resultado en base de datos u otra ubicación. \\
Results display & Presentación del resultado para su uso por el área operativa. \\
\bottomrule
\end{tabularx}
\end{table*}

\section*{4. Desarrollo del tema}

\cajaitem{Idea central}{Una aplicación Big Data de calidad surge cuando se unen dos dimensiones: capacidades técnicas de desarrollo y comprensión del servicio que se desea resolver. A partir de esa base, el trabajo se organiza en un proceso secuencial que transforma una necesidad en resultados utilizables.}

\subsection*{4.1. ¿Qué hace posible una buena aplicación Big Data?}

El material responde esta pregunta a partir de dos grupos de habilidades:

\begin{enumerate}[leftmargin=*]
    \item \textbf{Fundamentos de programación}
    \begin{itemize}[leftmargin=*]
        \item Capacidad de programar con \textbf{Java} y \textbf{Scala}.
        \item Familiaridad con \textbf{SQL}.
        \item Manejo de \textbf{operaciones en Linux}.
    \end{itemize}
    
    \item \textbf{Experiencia en desarrollo de servicios}
    \begin{itemize}[leftmargin=*]
        \item Comprensión de los \textbf{procesos de desarrollo}.
        \item Comprensión del \textbf{contexto o trasfondo del servicio}.
    \end{itemize}
\end{enumerate}

Esto implica que el desarrollador Big Data no solo debe saber escribir código, sino también comprender \textbf{qué problema se está resolviendo}, \textbf{cómo se integra la solución al proceso organizacional} y \textbf{cómo se traduce una necesidad del negocio en resultados útiles}.

\subsection*{4.2. Relación entre conocimiento técnico y conocimiento del servicio}

Según el esquema del material, ambos componentes son complementarios:

\begin{itemize}[leftmargin=*]
    \item El \textbf{conocimiento técnico} permite construir la solución.
    \item El \textbf{conocimiento del servicio} permite orientar la solución hacia un objetivo real.
\end{itemize}

Si falta la primera parte, no se puede implementar la aplicación. Si falta la segunda, la aplicación puede funcionar técnicamente, pero no responder al problema correcto.

\subsection*{4.3. Desarrollo como proceso y no como acto aislado}

El documento no presenta el desarrollo como una sola fase de programación, sino como un \textbf{flujo completo}. La programación aparece en la etapa de \textit{data analysis}, pero antes de eso ya hubo trabajo de requerimientos, selección tecnológica y validación; y después de programar todavía faltan almacenamiento y presentación de resultados.

Esto enseña una idea importante para evaluación: \textbf{desarrollar una aplicación Big Data es recorrer un ciclo completo}, no solamente escribir consultas o código.

\cajaitem{Errores comunes o confusiones frecuentes}{
\begin{itemize}[leftmargin=*]
    \item Creer que desarrollar una aplicación Big Data es solamente programar.
    \item Pensar que la selección tecnológica debe hacerse usando la tecnología más nueva sin evaluar estabilidad o viabilidad.
    \item Omitir el análisis de requerimientos y pasar directamente a la implementación.
    \item Confundir requerimientos del negocio con indicadores medibles.
    \item Olvidar que los resultados deben almacenarse y mostrarse para su uso operativo.
\end{itemize}
}

\subsection*{4.4. Qué enfatiza realmente esta sección}

Esta parte del capítulo es \textbf{breve y orientada al proceso}. No entra en detalles extensos de algoritmos ni de implementación de código. Su foco está en:

\begin{itemize}[leftmargin=*]
    \item las competencias mínimas del desarrollador,
    \item la secuencia general de desarrollo,
    \item y los puntos que deben cuidarse antes, durante y después del análisis de datos.
\end{itemize}

Por ello, para un cuestionario, lo más probable es que se pregunten \textbf{etapas}, \textbf{relaciones entre habilidades} y \textbf{finalidad de cada fase}.

\section*{5. Procesos, arquitectura o flujo del tema}

En esta sección, el PDF enfatiza más un \textbf{flujo de desarrollo} que una arquitectura técnica detallada. El proceso general es el siguiente:

\begin{enumerate}[leftmargin=*]
    \item \textbf{Requirement submission}  
    Se inicia con negociación entre múltiples departamentos y un análisis preliminar de factibilidad del proyecto.
    
    \item \textbf{Requirement analysis}  
    Se realiza levantamiento de requerimientos y estudio de mercado.
    
    \item \textbf{Technology selection}  
    Se define el diagrama de flujo de datos y la solución técnica.
    
    \item \textbf{Feasibility analysis}  
    Se construye la plataforma técnica y se hacen pruebas de factibilidad.
    
    \item \textbf{Indicator analysis}  
    Se convierten los requerimientos en indicadores.
    
    \item \textbf{Data interaction}  
    Incluye recolección de datos, limpieza, almacenamiento de datos fuente y middleware.
    
    \item \textbf{Data analysis}  
    Se implementa el código para procesar y analizar los datos.
    
    \item \textbf{Result storage}  
    Los resultados se guardan en una base de datos u otra ubicación.
    
    \item \textbf{Results display}  
    Los resultados se presentan para uso del área de operaciones.
\end{enumerate}

\subsection*{5.1. Lectura del flujo como cadena lógica}

\begin{center}
\textbf{Necesidad} $\rightarrow$ \textbf{Análisis} $\rightarrow$ \textbf{Selección técnica} $\rightarrow$ \textbf{Validación} $\rightarrow$ \textbf{Preparación de datos} $\rightarrow$ \textbf{Implementación} $\rightarrow$ \textbf{Almacenamiento} $\rightarrow$ \textbf{Visualización}
\end{center}

\subsection*{5.2. Tabla resumen del proceso}

\begin{table*}[t]
\small
\centering
\begin{tabularx}{\textwidth}{>{\raggedright\arraybackslash}p{0.2\textwidth} >{\raggedright\arraybackslash}p{0.2\textwidth} >{\raggedright\arraybackslash}X}
\toprule
\textbf{Etapa} & \textbf{Propósito} & \textbf{Elemento clave mencionado en el PDF} \\
\midrule
Requirement submission & Acordar el inicio del proyecto & Negociación interdepartamental y factibilidad \\
Requirement analysis & Entender la necesidad & Encuesta de requerimientos y estudio de mercado \\
Technology selection & Definir cómo se hará & Diagrama de flujo de datos y solución técnica \\
Feasibility analysis & Verificar si es viable & Construcción de plataforma y pruebas \\
Indicator analysis & Traducir la necesidad a métricas & Conversión entre requerimientos e indicadores \\
Data interaction & Preparar el insumo de datos & Recolección, limpieza, almacenamiento fuente y middleware \\
Data analysis & Ejecutar la lógica de análisis & Implementación de código \\
Result storage & Conservar resultados & Base de datos u otra ubicación \\
Results display & Entregar valor al usuario final & Uso por el área operativa \\
\bottomrule
\end{tabularx}
\end{table*}

\section*{6. Herramientas, algoritmos o componentes importantes}

Aunque esta sección no desarrolla algoritmos específicos, sí menciona componentes y habilidades importantes para el desarrollo de aplicaciones Big Data.

\begin{table*}[t]
\small
\centering
\begin{tabularx}{\textwidth}{>{\raggedright\arraybackslash}p{0.22\textwidth} >{\raggedright\arraybackslash}p{0.24\textwidth} >{\raggedright\arraybackslash}X}
\toprule
\textbf{Herramienta o componente} & \textbf{¿Qué es o qué representa?} & \textbf{¿Para qué sirve dentro del tema?} \\
\midrule
Java & Lenguaje de programación & Base para implementar aplicaciones Big Data \\
Scala & Lenguaje de programación & Base técnica importante para desarrollar soluciones \\
SQL & Lenguaje de consulta & Permite trabajar con datos de forma estructurada \\
Linux & Entorno operativo & Soporta la operación técnica del desarrollo \\
Diagrama de flujo de datos & Representación del movimiento de los datos & Ayuda en la selección tecnológica y diseño de la solución \\
Plataforma técnica & Entorno construido para el proyecto & Permite probar factibilidad antes de producción \\
Middleware & Capa intermedia del sistema & Participa en la interacción y manejo de datos \\
Base de datos u otra ubicación & Destino del resultado & Permite persistir la salida del análisis \\
\bottomrule
\end{tabularx}
\end{table*}

\subsection*{6.1. Lo que conviene recordar}

\begin{itemize}[leftmargin=*]
    \item En esta sección, \textbf{Java, Scala, SQL y Linux} aparecen como base de competencias.
    \item El contenido \textbf{no se centra en algoritmos concretos}, sino en el proceso de desarrollo.
    \item El \textbf{diagrama de flujo de datos} y la \textbf{plataforma técnica} son piezas clave de diseño y validación.
\end{itemize}

\section*{7. Aplicaciones o ejemplos}

El material de esta parte no desarrolla un caso completo de implementación, pero sí permite reconstruir aplicaciones y usos probables del proceso.

\subsection*{7.1. Ejemplo general de aplicación}

Una organización detecta una necesidad y presenta un requerimiento. Luego:

\begin{itemize}[leftmargin=*]
    \item varias áreas negocian la viabilidad del proyecto,
    \item se analiza la necesidad y el mercado,
    \item se selecciona una tecnología adecuada,
    \item se construye la plataforma y se prueba,
    \item se recolectan y limpian los datos,
    \item se implementa el análisis,
    \item los resultados se guardan,
    \item y finalmente se muestran al área operativa.
\end{itemize}

Este ejemplo refleja exactamente la lógica del flujo presentado en el documento.

\subsection*{7.2. Relación con aplicaciones Big Data del capítulo}

Dentro del mismo capítulo, Big Data aparece aplicado en dominios como:

\begin{itemize}[leftmargin=*]
    \item \textbf{Finanzas}: transacciones de alta frecuencia y análisis de riesgo crediticio.
    \item \textbf{Logística}: logística inteligente, control de costos y gestión de riesgos.
    \item \textbf{Salud}: predicción de epidemias y gestión sanitaria inteligente.
    \item \textbf{Internet}: perfiles de usuario, recomendación personalizada y publicidad.
    \item \textbf{Ciudad}: transporte inteligente y planificación urbana.
\end{itemize}

Estos ejemplos ayudan a entender que el proceso de desarrollo estudiado no se limita a un sector, sino que puede aplicarse en múltiples escenarios.

\subsection*{7.3. Precisión importante}

El PDF no detalla aquí un ejemplo de código ni una arquitectura de despliegue completa para esta sección específica. Por eso, el estudio debe concentrarse en \textbf{la lógica del desarrollo}, \textbf{las etapas} y \textbf{las competencias requeridas}.

\section*{8. Resumen para memorizar}

\begin{itemize}[leftmargin=*]
    \item Una buena aplicación Big Data requiere \textbf{habilidades técnicas} y \textbf{comprensión del servicio}.
    \item Las bases técnicas mencionadas son \textbf{Java, Scala, SQL y Linux}.
    \item También se necesita entender \textbf{el proceso de desarrollo} y \textbf{el contexto del negocio}.
    \item El desarrollo sigue una secuencia:
    \begin{center}
        \textbf{requerimiento $\rightarrow$ análisis $\rightarrow$ selección tecnológica $\rightarrow$ factibilidad $\rightarrow$ indicadores $\rightarrow$ interacción de datos $\rightarrow$ análisis $\rightarrow$ almacenamiento $\rightarrow$ visualización}
    \end{center}
    \item \textbf{Technology selection} incluye el \textbf{diagrama de flujo de datos} y la \textbf{solución técnica}.
    \item \textbf{Feasibility analysis} consiste en construir la plataforma técnica y validar su funcionamiento.
    \item \textbf{Data interaction} incluye \textbf{recolección, limpieza, almacenamiento fuente y middleware}.
    \item \textbf{Data analysis} corresponde a la \textbf{implementación del código}.
    \item \textbf{Result storage} guarda la salida en base de datos u otra ubicación.
    \item \textbf{Results display} entrega el resultado al área operativa.
    \item No siempre debe elegirse la tecnología más nueva; primero debe evaluarse su \textbf{viabilidad y estabilidad}.
\end{itemize}

\section*{9. Preguntas probables de cuestionario}

\begin{enumerate}[leftmargin=*]
    \item \textbf{¿Qué dos grupos de capacidades hacen posible una excelente aplicación Big Data?}\\
    \textbf{Respuesta corta:} fundamentos de programación y experiencia en desarrollo de servicios.
    
    \item \textbf{¿Qué lenguajes de programación menciona el material como base?}\\
    \textbf{Respuesta corta:} Java y Scala.
    
    \item \textbf{¿Qué lenguaje de consulta debe conocer el desarrollador?}\\
    \textbf{Respuesta corta:} SQL.
    
    \item \textbf{¿Qué entorno operativo se menciona como conocimiento necesario?}\\
    \textbf{Respuesta corta:} operaciones en Linux.
    
    \item \textbf{¿Qué incluye la etapa de \textit{requirement submission}?}\\
    \textbf{Respuesta corta:} negociación entre múltiples departamentos y análisis de factibilidad del proyecto.
    
    \item \textbf{¿Qué actividades forman parte de \textit{requirement analysis}?}\\
    \textbf{Respuesta corta:} levantamiento de requerimientos y estudio de mercado.
    
    \item \textbf{¿Qué se define en la etapa de \textit{technology selection}?}\\
    \textbf{Respuesta corta:} el diagrama de flujo de datos y la solución técnica.
    
    \item \textbf{¿Para qué sirve \textit{feasibility analysis}?}\\
    \textbf{Respuesta corta:} para construir la plataforma técnica y probar si la solución es viable.
    
    \item \textbf{¿Qué significa \textit{indicator analysis}?}\\
    \textbf{Respuesta corta:} convertir los requerimientos en indicadores.
    
    \item \textbf{¿Qué tareas incluye \textit{data interaction}?}\\
    \textbf{Respuesta corta:} recolección de datos, limpieza, almacenamiento de datos fuente y middleware.
    
    \item \textbf{¿Qué ocurre en \textit{data analysis}?}\\
    \textbf{Respuesta corta:} se implementa el código para analizar o procesar los datos.
    
    \item \textbf{¿Dónde se guardan los resultados del análisis?}\\
    \textbf{Respuesta corta:} en una base de datos u otra ubicación.
    
    \item \textbf{¿Quién usa la salida de la etapa \textit{results display}?}\\
    \textbf{Respuesta corta:} el área de operaciones.
    
    \item \textbf{¿Es correcto elegir siempre la tecnología más nueva sin considerar estabilidad?}\\
    \textbf{Respuesta corta:} no; la selección debe considerar viabilidad y estabilidad.
    
    \item \textbf{¿Puede usarse un host en la nube o una máquina física como servidor?}\\
    \textbf{Respuesta corta:} sí, el material lo deja como posibilidad en la evaluación tipo cuestionario.
\end{enumerate}

\section*{10. Mini glosario}

\begin{description}
    \item[Big Data Application Development] Desarrollo de aplicaciones orientadas al manejo y aprovechamiento de grandes volúmenes de datos.
    \item[Programming Basics] Conocimientos técnicos base requeridos para implementar la solución.
    \item[Service Development Experience] Comprensión del proceso de desarrollo y del contexto del servicio.
    \item[Requirement Submission] Presentación inicial del requerimiento con negociación y revisión preliminar de factibilidad.
    \item[Requirement Analysis] Estudio detallado de necesidades y del entorno o mercado.
    \item[Technology Selection] Elección de la solución técnica adecuada y definición del flujo de datos.
    \item[Feasibility Analysis] Validación práctica de que la solución puede construirse y funcionar.
    \item[Indicator Analysis] Traducción de necesidades del negocio a indicadores observables o medibles.
    \item[Data Interaction] Manejo inicial de los datos: recolectar, limpiar, almacenar y conectar componentes intermedios.
    \item[Data Analysis] Etapa en la que se implementa la lógica analítica mediante código.
    \item[Result Storage] Persistencia de los resultados obtenidos.
    \item[Results Display] Presentación final de los resultados para su uso operativo.
    \item[Middleware] Componente intermedio que apoya la interacción y circulación de datos entre partes del sistema.
\end{description}

\end{document}
